% ===================================================================
%  QESPM Supplementary Material
%  Supporting figures for the main manuscript
%  Compile with: pdflatex supplementary
% ===================================================================

\documentclass[11pt,a4paper]{article}
\usepackage[utf8]{inputenc}
\usepackage[T1]{fontenc}
\usepackage{textcomp}
\usepackage{amsmath,amssymb}
\usepackage{graphicx}
\usepackage{hyperref}
\usepackage[margin=2.5cm]{geometry}
\renewcommand{\thesection}{S\arabic{section}}
\renewcommand{\thefigure}{S\arabic{figure}}
\renewcommand{\thetable}{S\arabic{table}}
\renewcommand{\theequation}{S\arabic{equation}}

\title{Supplementary Material:\\
       Quantum Limits of Surface Metrology:\\
       Spatial Information Extraction by a Single Trapped Ion}

\author{Dawid Kucharski\\
        Poznan University of Technology, Pozna\'n, Poland}

\date{\today}

\begin{document}
\maketitle

\noindent
\textbf{Supplementary file description.} Full derivations, validation
details, robustness tests, Monte Carlo checks, calibration and traceability
chain, symbol table, experimental roadmap and forward-model derivations
supporting the main manuscript.

Sections S2--S16 contain the full derivations and validation details that
support the main manuscript; all material is cross-referenced therein as
Fig.~S1--S16 and Tables~S1--S2.

% ===================================================================
\section{Supplementary Figures}
% ===================================================================

% ---- Scanning geometry: shown as Fig.~1 of the main text, not duplicated ----
\noindent
\textbf{Scanning geometry.}
The QESPM scanning geometry and protocol are shown in Fig.~1 of the main
text: the ion is held stationary in a planar trap while the sample is
translated laterally beneath it, and the secular frequency shift
$\Delta\omega_x(x,y)$ is recorded at each position.  The ion--surface
separation $h$ is maintained by adjusting the trap-electrode voltages.
The scan range is limited by the field of view of the ion trap
($64\,\mu\text{m}$ in the simulation).

% ---- Figure S2: Signal map ----
\begin{figure}[htbp]
\centering
\includegraphics[width=0.85\textwidth]{figS4_signal_map.pdf}
\caption{\textbf{Simulated ion signal map.}
Peak signal amplitude $|\Delta\omega_x|/2\pi$ for a sinusoidal surface of
amplitude $A = 100$\,nm as a function of surface wavelength $\lambda$ and
ion height $h$, computed from the analytical ITF.  White--magenta contours
mark the 5\,Hz--50\,kHz signal levels (the 5\,Hz contour is the SNR$=1$
limit for the conservative noise floor).  The white dashed line is the
optimal wavelength $\lambda_{\rm opt} = \pi h$; the star marks the nominal
operating point ($\lambda = 40\,\mu\text{m}$, $h = 40\,\mu\text{m}$).}
\label{fig:S_signal}
\end{figure}

% ---- Figure S3: Micromotion analysis ----
\begin{figure}[htbp]
\centering
\includegraphics[width=0.85\textwidth]{fig8_micromotion.pdf}
\caption{\textbf{Micromotion as a complementary observable.}
Comparison of the two surface-sensing channels at $h = 40\,\mu\text{m}$.
(a) Normalised ITF magnitude for the secular channel
$\Delta\omega_x$ ($\propto k^2 e^{-kh}$, curvature response) and the
excess-micromotion channel $\beta_x$ ($\propto |k| e^{-kh}$, gradient
response).  The micromotion channel peaks at $k = 1/h$
($\lambda = 2\pi h$) and is relatively more sensitive at lower spatial
frequencies than the secular channel, which peaks at $k = 2/h$.
(b) Complementary spatial-frequency coverage of the two channels: the
gradient and curvature responses together span the full accessible
$k$-band.  Because the two channels scale differently with $h$ and $k$,
their ratio provides an independent diagnostic (see the main text).}
\label{fig:S_micromotion}
\end{figure}

% ---- Figure S4: Degeneracy theorem ----
\begin{figure}[htbp]
\centering
\includegraphics[width=0.85\textwidth]{fig9_degeneracy.pdf}
\caption{\textbf{Topography--charge identifiability limit --- numerical illustration.}
Ratio of the frequency-shift vectors at two ion heights,
$\Delta\omega(h_1)/\Delta\omega(h_2) = e^{-k(h_1-h_2)}$ for
$h_1 = 20\,\mu\text{m}$ and $h_2 = 80\,\mu\text{m}$.  The ratio depends
only on $k$ and the heights --- not on the $(z_s, \sigma)$ mixture ---
because both the topographic and the charge contributions propagate
identically as $e^{-kh}$.  The two-height forward matrix is therefore
rank~1, and \textbf{no finite number of heights}, without independent
information, can separate the two contributions (Theorem in the main
text).  Resolution strategies are indicated:
in-situ charge elimination, a known reference scan, or an independent
material-contrast channel such as the motional heating rate.}
\label{fig:S_degeneracy}
\end{figure}

% ---- Figure S5: AFM input data ----
\begin{figure}[htbp]
\centering
\includegraphics[width=0.85\textwidth]{figV1_afm_input.pdf}
\caption{\textbf{AFM input topography --- stainless steel.}
Input surface topography used for the real-data validation
(Sec.~4.2 of the main text).  The data are from
Camargo~Jr.\ (2019, Mendeley Data,
\url{https://doi.org/10.17632/6dzmrjngcg.3},
CC~BY~4.0).  (a) The $20 \times 20\,\mu\text{m}^2$ AFM scan of
stainless steel.  (b) Power spectral density (PSD) showing roughness
power concentrated at wavelengths $\lambda \sim 2$--$20\,\mu\text{m}$,
which overlap with the QESPM ITF passband for $h \lesssim 10\,\mu\text{m}$.}
\label{fig:S_afm_input}
\end{figure}

% ---- Figure S6: Ion signal from AFM data ----
\begin{figure}[htbp]
\centering
\includegraphics[width=0.85\textwidth]{figV2_afm_ion_signal.pdf}
\caption{\textbf{Simulated ion signal from AFM topography.}
Forward-model prediction of $\Delta\omega_x(x,y)$ computed from the
AFM-measured stainless steel topography at four ion heights:
$h = 2, 5, 10, 40\,\mu\text{m}$.  The signal amplitude grows by
$\approx 2.8\times10^5$ from $h = 40\,\mu\text{m}$ to $h = 2\,\mu\text{m}$,
consistent with the $e^{-kh}$ height scaling of the ITF.  At
$h \lesssim 5\,\mu\text{m}$ the predicted shifts reach $|\Delta\omega_x|/2\pi
\sim 10^{14}$--$10^{15}$\,Hz, comparable to or exceeding $\omega_x$ itself:
these panels illustrate the asymptotic ITF behaviour outside the
perturbative regime, not an experimentally realisable configuration
(cf.~the heating-rate caveat in the main text).  Note the progressive
emergence of fine spatial features as $h$ decreases.}
\label{fig:S_afm_signal}
\end{figure}

% ---- Figure S7: AFM reconstruction ----
\begin{figure}[htbp]
\centering
\includegraphics[width=0.85\textwidth]{figV3_afm_reconstruction.pdf}
\caption{\textbf{QESPM reconstruction from AFM data.}
Tikhonov-regularised reconstruction of the stainless steel surface from
the simulated ion signal at $h = 40, 10, 5, 2\,\mu\text{m}$.  (a)--(d) Reconstructed topography.  (e)--(h) Residual
(reconstruction $-$ true AFM topography).  The Pearson correlation
improves from $r = 0.62$ at $h = 40\,\mu\text{m}$ to $r = 0.78$ at
$h = 2\,\mu\text{m}$, confirming the predicted height dependence of the
reconstruction quality.}
\label{fig:S_afm_recon}
\end{figure}

% ===================================================================
\section{Floquet analysis: beyond the pseudopotential approximation}
% ===================================================================

The pseudopotential approximation used in the main text is valid for
$|a|, q^2 \ll 1$.  This section verifies that the full Floquet solution of
the Mathieu equation reproduces the pseudopotential result to
$\mathcal{O}(q^2)$ and computes the leading correction.

The exact 1D equation of motion for an ion in a Paul trap with a static
surface perturbation $F_s(x) = -e\,\partial\Phi/\partial x$ is, in
dimensionless form ($\tau = \Omega t/2$, $u = x/R_{\rm rf}$, where
$R_{\rm rf}$ is the characteristic trap dimension --- the distance
from the trap centre to the nearest rf electrode):
\begin{equation}
\frac{d^2 u}{d\tau^2} + [a - 2q\cos(2\tau)]u
    = -\frac{4e}{m\Omega^2 R_{\rm rf}}
       \left.\frac{\partial\Phi}{\partial x}\right|_{x = R_{\rm rf}u}.
\label{eq:mathieu_perturbed}
\end{equation}
For a sinusoidal surface with $|A| \ll h$, the perturbation is
$F_s \propto \cos(k_s x)$, which couples the secular motion to the rf drive
through the position dependence of the force.  Expanding about the ion's
equilibrium position $u_0$ and retaining only the linear term (valid for
small oscillation amplitudes), Eq.~(\ref{eq:mathieu_perturbed}) takes the
form of a perturbed Mathieu equation:
\begin{equation}
\frac{d^2\delta u}{d\tau^2} + [a + \delta a - 2q\cos(2\tau)]\delta u = 0,
\label{eq:mathieu_pert_linear}
\end{equation}
with the \textbf{effective Mathieu parameter shift}
\begin{equation}
\delta a = \frac{8e}{m\Omega^2 R_{\rm rf}^2}
        \left.\frac{\partial^2\Phi}{\partial x^2}\right|_{x_0}
      = -\frac{8e E_{\rm bg}}{m\Omega^2 R_{\rm rf}^2}\,
        A\,k_s^2\,\sin(k_s x_0)\,e^{-k_s h}.
\label{eq:delta_a}
\end{equation}

The Floquet characteristic exponent $\mu(a,q)$ for the unperturbed Mathieu
equation satisfies $\cosh[\mu(a,q)] = 1 - 2\Delta(a,q)$, where $\Delta(a,q)$
is the Hill determinant.  For $|a|, q^2 \ll 1$, the expansion
$\mu(a,q) \approx i\sqrt{a + q^2/2} + \mathcal{O}(a^2, q^4)$ recovers the
pseudopotential secular frequency $\omega_{\rm sec} = (\Omega/2)\,{\rm
Im}[\mu]$.  The surface perturbation shifts $a \to a + \delta a$, and the
frequency shift to first order in $\delta a$ is
\begin{equation}
\Delta\omega_x^{\rm(Floquet)} = \frac{\Omega}{2}\,
    \left.\frac{\partial\,{\rm Im}[\mu]}{\partial a}\right|_{a,q}\delta a
    = \frac{e}{2m\omega_x}\,
       \left.\frac{\partial^2\Phi}{\partial x^2}\right|_{x_0}
       \left[1 + \mathcal{O}(q^2)\right].
\label{eq:floquet_result}
\end{equation}

The $\mathcal{O}(q^2)$ correction can be evaluated numerically from the
power-series expansion of the Mathieu characteristic exponent
$\beta(a,q)$ (McLachlan 1947).  For $q = 0.3$, the exact Floquet
result differs from the pseudopotential by $\approx 3\%$; for $q = 0.5$, the
difference reaches $\approx 11\%$.  Table~\ref{tab:floquet} summarises the
correction, computed from the Mathieu characteristic exponent via the Hill
determinant (equivalently, the monodromy matrix of the Mathieu equation
over one rf period, $\beta=(1/\pi)\arccos(\operatorname{tr} M/2)$).  The
pseudopotential formula of the main text is therefore accurate to better
than $11\%$ for all practical trap configurations, and the exact Floquet
result can be substituted when higher precision is required.  Throughout
the main text, the pseudopotential expressions are used with the
understanding that the $\mathcal{O}(q^2)$ correction is a known systematic
that can be calibrated out.

\begin{table}[htb]
\centering
\caption{Exact Floquet versus pseudopotential response: relative excess of
the exact frequency-response derivative $\partial\,\mathrm{Im}\,\mu/\partial a$
over its pseudopotential value, evaluated at $a=0$ for the $\pi$-periodic
Mathieu equation $u''+[a-2q\cos(2\tau)]u=0$.  The text figures of
${\approx}3\%$ at $q=0.3$ and ${\approx}11\%$ at $q=0.5$ are reproduced.}
\label{tab:floquet}
\small
\begin{tabular}{@{}lccccccc@{}}
\hline
$q$ & 0.1 & 0.2 & 0.3 & 0.4 & 0.5 & 0.6 & 0.7 \\
\hline
correction & $0.31\%$ & $1.28\%$ & $3.07\%$ & $6.00\%$ & $10.74\%$ & $18.71\%$ & $33.68\%$ \\
\hline
\end{tabular}
\end{table}

\textbf{Micromotion coupling to surface topography.}
Even in a well-compensated trap, the ion executes driven micromotion
$x_\mu(t) = -(q/2)x_{\rm sec}\cos(\Omega t)$, where $x_{\rm sec}$ is the
secular displacement.  The surface force $F_s(x_0 + x_\mu(t))$ sampled at
the rf frequency produces sidebands in the motional spectrum at
$\omega_x \pm \Omega$.  Expanding $F_s$ to first order in $x_\mu$:
\begin{equation}
F_s(x_0 + x_\mu) = F_s(x_0) - \frac{q}{2}x_{\rm sec}
    \left.\frac{\partial F_s}{\partial x}\right|_{x_0}\!\cos(\Omega t)
    + \mathcal{O}(x_\mu^2).
\end{equation}
The sideband amplitude relative to the carrier is
$|S_{\pm}|/|S_0| \approx (q/2)(x_{\rm sec}/\lambda_{\rm feature})$,
where $\lambda_{\rm feature} \sim 2\pi/k$ is the surface feature scale.
For $q = 0.3$, $x_{\rm sec} \sim 10\,\textup{nm}$ (ground-state-cooled ion),
and $\lambda_{\rm feature} \sim 30\,\mu\textup{m}$, the sideband ratio is
${\sim}5 \times 10^{-7}$, below the measurement noise floor.  However,
for warmer ions ($\bar{n} \sim 10^3$, $x_{\rm sec} \sim 300\,\textup{nm}$)
or at small $h$ where $\lambda_{\rm feature}$ is smaller, the sidebands
become detectable and provide an independent diagnostic channel.

\textbf{Sideband-channel transfer function.}
The sideband channel carries the gradient transfer function
$H_{\rm mm}(k;h)\propto k_x e^{-kh}$, whose magnitude peaks at $k=1/h$
($\lambda=2\pi h$) --- one octave below the curvature channel's peak at
$k=2/h$ --- so the two channels jointly span the accessible band
(Fig.~S2).  Because $H_{\rm mm}$ is a linear functional of
the same scalar field $\varphi_s$ of the topography--charge degeneracy
corollary of the main text, the sideband channel cannot break the
topography--charge degeneracy; its role is as an independent consistency
check on $\varphi_s$ and as coverage of the low-$k$ wing of the band,
where the curvature channel is insensitive.

\textbf{Carrier (rf-averaged) correction.}
The secular frequency is measured on the rf-time-averaged orbit, so the
carrier shift samples the curvature convoluted with the micromotion
distribution.  Taylor-expanding
$\partial_x^2\Phi(x_0+x_\mu(t))$ and averaging over the rf cycle,
$\langle x_\mu^2\rangle = q^2x_{\rm sec}^2/8$, gives
\begin{equation}
\Delta\omega_x^{\rm(carrier)} = \Delta\omega_x
    \Bigl[1 + \frac{q^2\pi^2}{4}\Bigl(\frac{x_{\rm sec}}{\lambda_{\rm
feature}}\Bigr)^{2} + \ldots\Bigr],
\end{equation}
since $\partial_x^4\Phi/\partial_x^2\Phi = (2\pi/\lambda_{\rm feature})^2$
for a sinusoidal feature.  For $q=0.3$ and a cold ion
($x_{\rm sec}\sim10$\,nm) at $\lambda_{\rm feature}\sim30$\,\textmu m
the correction is ${\sim}2\times10^{-8}$ --- negligible --- but for a warm
ion ($x_{\rm sec}\sim300$\,nm) at micrometre-scale features it reaches
${\sim}2\%$, and must be entered into the Type-B uncertainty budget as a
bounded systematic with a known analytic form.

\textbf{Anharmonic corrections.}
The surface perturbation contributes terms beyond quadratic in the ion
displacement:
$\Phi^{(1)}(x) = \Phi^{(1)}(x_0) + \partial_x\Phi^{(1)}x
+ \tfrac{1}{2}\partial_x^2\Phi^{(1)}x^2
+ \tfrac{1}{6}\partial_x^3\Phi^{(1)}x^3
+ \tfrac{1}{24}\partial_x^4\Phi^{(1)}x^4 + \cdots$.
The quartic term shifts the secular frequency at $\mathcal{O}(x_{\rm zp}^2)$
via first-order perturbation theory in the anharmonicity: normal-ordering
$(a+a^\dagger)^4 = 6(a^\dagger a)^2 + 6a^\dagger a + 3$ makes the
$|n\rangle\to|n{+}1\rangle$ level spacing shift by
$(e/4m\omega_x)\,\partial_x^4\Phi^{(1)}\,x_{\rm zp}^2$, i.e.
\begin{equation}
\delta\omega_x^{\rm(anh)} = \frac{e}{4m\omega_x}\,
    \left.\frac{\partial^4\Phi^{(1)}}{\partial x^4}\right|_{x_0}\!
    x_{\rm zp}^2,
\end{equation}
(the odd cubic term $\partial_x^3\Phi^{(1)}x^3/6$ produces no first-order
frequency shift).  Here $x_{\rm zp} = \sqrt{\hbar/2m\omega_x} \approx
13\,\textup{nm}$ for ${}^{40}$Ca$^+$ at $\omega_x/2\pi = 1\,\textup{MHz}$.
For a sinusoidal surface, $\partial_x^4\Phi^{(1)} \propto k^4\Phi^{(1)}$,
giving
$\delta\omega_x^{\rm(anh)}/\Delta\omega_x \sim \tfrac12 x_{\rm zp}^2 k^2$.
At $h = 40\,\mu\textup{m}$ and $k = 2 \times 10^5\,\textup{rad/m}$,
this ratio is ${\sim}3 \times 10^{-6}$ --- completely negligible.
At $h = 5\,\mu\textup{m}$ the band edge $kh = 13$ gives
${\sim}6 \times 10^{-4}$, still below the $5\,\textup{Hz}$ noise floor for
the nominal surface amplitude.  Anharmonicity becomes the dominant
correction only for $h \lesssim 1\,\mu\textup{m}$, where Casimir--Polder
forces already invalidate the purely electrostatic model.

\begin{figure}[htb]
\centering
\includegraphics[width=\textwidth]{figS1_trajectory.pdf}
\caption{Direct numerical integration of the ion's secular equation of
motion near a sinusoidal surface ($A = 200$\,nm, $\lambda = 40$\,\textmu m,
$h = 40$\,\textmu m, trap centre at $x_{\rm scan} = \lambda/4$).
(a) Trajectory over $5$\,ms integration.
(b) Zoom on $15$ secular periods with surface (blue) and without
(red dashed).  (c) Power spectral density with the ITF-predicted
secular peak (red dashed) compared to the unperturbed reference (grey
dotted).  (d) Zoom on the secular peak: the ITF predicts
$|\Delta f| \approx 2.8$\,kHz ($\Delta\omega/\omega \approx 0.28\%$),
consistent with the analytical model.  The ion motion remains stable
and harmonic, validating the perturbative treatment.}
\label{fig:trajectory}
\end{figure}

Direct numerical integration confirms that the ion's secular motion
remains stable and well-described by the linear perturbation theory for
topographic amplitudes up to $200$\,nm, providing numerical validation of
the analytical ITF.

% ===================================================================
\section{Exact (Rayleigh--Floquet) cross-validation}
% ===================================================================

To validate the first-order boundary perturbation expansion against an exact
electrostatic solution, the Laplace equation with the corrugated Dirichlet
boundary is solved exactly by the Rayleigh (Floquet) mode expansion
$\Phi(x,z) = -E_{\rm bg}z + \sum_n c_n \sin(nkx)e^{-nkz}$, with the
coefficients $c_n$ determined by the exact boundary condition
$\Phi(x, A\sin(kx)) = 0$.
Figure~\ref{fig:bem}(a) compares the secular frequency shift computed
exactly and via the linear ITF for a sinusoidal surface of period
$\lambda = 40$\,\textmu m as a function of amplitude $A$.
The ITF reproduces the exact result at small $A/h$ (relative error
$0.01\%$ at $A = 100$\,nm), with deviations growing systematically as
$0.14\,(kA)^2$ at larger amplitudes
($1.4\%$ at $A = 2$\,\textmu m) as the nonlinear terms
$\mathcal{O}(z_s^2)$ become significant (Fig.~\ref{fig:bem}b).
\textbf{Convergence of the exact solver.}
The Rayleigh (Floquet) expansion is known to converge for grating
amplitudes $kA \lesssim 0.448$ (Petit 1966; Millar 1973);
the entire amplitude range of Fig.~\ref{fig:bem} ($kA \le 0.31$) lies
within this domain.  With the complete basis ($N=60$ sin and cos
harmonics each), the relative boundary
residual $\max_x|\Phi(x,A\sin kx)|/|E_{\rm bg}A|$ is
$7.9\times10^{-4}$ at $A=10$\,nm, $7.9\times10^{-3}$ at $A=100$\,nm and
$0.15$ at $A=2$\,\textmu m (the residual is carried by high harmonics
that decay as $e^{-nkh}$ toward the ion); the frequency-shift observable
is converged far more rapidly: $|\Delta\omega_x(N{=}120)-
\Delta\omega_x(N{=}60)|/|\Delta\omega_x|<10^{-8}$ at all plotted
amplitudes, so the quoted exact-vs-ITF errors are unaffected by
truncation at the $10^{-4}$ level.

\begin{figure}[htb]
\centering
\includegraphics[width=\textwidth]{figV7_rayleigh_validation.pdf}
\caption{Exact (Rayleigh) validation of the linear ITF.
(a) Frequency shift vs.\ surface amplitude: exact (Rayleigh) and ITF (linear).
(b) Relative error as a function of amplitude, scaling as $(kA)^2$.
(c) Relative error vs.\ $A/h$ ratio, showing increasing deviation
at larger amplitudes.}
\label{fig:bem}
\end{figure}

\begin{figure}[htb]
\centering
\includegraphics[width=0.9\textwidth]{figO6_bem.pdf}
\caption{Boundary-element cross-check (single-layer collocation with the
periodic Green's function).  (a) Secular frequency shift at the crest from
the BEM solution, the linear ITF, and the exact (complete-basis) Rayleigh
solution, vs.\ surface amplitude.  (b) Relative errors of the BEM and of
the linear ITF against the exact solution: the BEM agrees to $\le0.6\%$
across the amplitude range, independently confirming the ITF at the
$1\%$ level without any use of the ITF itself.}
\label{fig:bem_cross}
\end{figure}

As a further cross-check from an independent discretisation family, the
Laplace equation is also solved by finite differences (FD) on the
corrugated domain with a grounded boundary (Fig.~\ref{fig:fd}): the crest
frequency shift ($A = 100$\,nm, $h = 40$\,\textmu m) converges with the
number of vertical grid points $N_v = 60$--$200$ toward the ITF
prediction, confirming the first-order boundary-perturbation result with
three mutually independent solvers (exact Rayleigh, BEM, FD).

\begin{figure}[htb]
\centering
\includegraphics[width=\textwidth]{figI14_fem.pdf}
\caption{Finite-difference cross-check of the linear ITF.
(a) Convergence of the crest frequency shift ($A = 100$\,nm,
$h = 40$\,\textmu m) with the number of vertical grid points $N_v$
($60$--$200$), approaching the ITF prediction (relative error vs.\ ITF).
(b) The corrugated computational domain with grounded boundary, showing
the ion plane at $h = 40$\,\textmu m.}
\label{fig:fd}
\end{figure}

% ===================================================================
\section{Robustness to model violation}
% ===================================================================

Three key model assumptions are systematically violated to map the
breakdown of the reconstruction (Fig.~\ref{fig:violation}):
(i)~\textbf{Non-perfect conductor}: a $d=2$\,\textmu m dielectric overlayer
($\varepsilon_r = 2$--$10$) on the conducting surface modifies the surface
potential by the exact three-layer factor
$F(k)=\varepsilon_r e^{kd}/[\varepsilon_r\cosh(kd)+\sinh(kd)]$
(derived in the main text), which
enhances the apparent heights by $+9\%$ to $+40\%$ across the passband;
(ii)~\textbf{non-uniform $E_{\rm bg}$}: a gradient
across the field of view; (iii)~\textbf{large surface slope}: violating
the $Ak \ll 1$ condition.  For each scenario, data are generated with the violated model and inverted assuming the
nominal (un-violated) model.  The reconstruction error is shown in
Fig.~\ref{fig:violation} versus the violation parameter for each
scenario: the amplitude error $|1 - R_q^{\rm rec}/R_q|$ for the
dielectric overlayer, which primarily rescales the reconstructed
heights (to $R_q^{\rm rec}>R_q$, since $F>1$), and the correlation loss $1-r$ for the field-gradient and
large-slope violations, which distort the reconstructed shape.

\begin{figure}[htb]
\centering
\includegraphics[width=\textwidth]{figV5_model_violation.pdf}
\caption{Robustness of the reconstruction to model violations.
(a) Dielectric overlayer ($d=2$\,\textmu m): amplitude error
$|1 - R_q^{\rm rec}/R_q|$ vs.\ $\varepsilon_r$, computed with the exact
three-layer factor $F(k)$ (which takes values $1.09$--$1.40$ across the
passband for $\varepsilon_r=2$--$10$).
(b) Non-uniform $E_{\rm bg}$: correlation loss $1-r$ vs.\ gradient amplitude.
(c) Large-slope breakdown: correlation loss $1-r$ vs.\ $Ak$ parameter.}
\label{fig:violation}
\end{figure}

% ===================================================================
\section{Monte Carlo uncertainty propagation}
% ===================================================================

The GUM-compliant budget of the main text is a single-point
evaluation.  Here, a full Monte Carlo propagation (JCGM 101:2008,
Supplement~1) is performed: $N_{\rm MC} = 5 \times 10^3$ realisations
of the forward + inversion chain are computed, sampling all six input
quantities from their uncertainty distributions.
Figure~\ref{fig:mc}(a) shows the reconstructed cross-section with the
95\% coverage interval for a sinusoidal grating
($A = 100$\,nm, $\lambda = 31.4$\,\textmu m, i.e.~$k = 2\times 10^5$
rad/m, the GUM evaluation point) at $h = 40$\,\textmu m.

The combined standard uncertainty reaches $u_c(z) \approx 300$\,nm at the
grating crest (Fig.~\ref{fig:mc}b), reproducing the GUM single-point
budget ($319$\,nm); the field-averaged value is ${\sim}200$\,nm because
the charge sensitivity is modulated sinusoidally across the field of view.
The distribution at the crest (Fig.~\ref{fig:mc}c) is approximately
Gaussian, confirming that the uncertainty propagation is well-behaved
despite the nonlinearity of the inversion.  The Monte Carlo result is thus
consistent with the GUM budget, the agreement confirming that the two
approaches (single-point linearisation and full nonlinear propagation)
describe the same dominant surface-charge uncertainty.  An explicit
coverage check closes the metrological loop: over $2\times10^{4}$ Monte
Carlo realisations at the crest, the fraction of true heights falling
inside the propagated $95\%$ interval ($\pm1.96\,u_c$ with
$u_c=319$\,nm) is $95.1\%$, with an estimate bias of $0.7$\,nm --- the
interval is correctly calibrated.

\begin{figure}[htb]
\centering
\includegraphics[width=\textwidth]{figV6_monte_carlo.pdf}
\caption{Monte Carlo uncertainty propagation.
(a) Reconstructed cross-section with 95\% coverage interval.
(b) Spatial map of combined standard uncertainty $u_c(z)$.
(c) Distribution of reconstructed height at the signal peak,
with 95\% coverage interval.}
\label{fig:mc}
\end{figure}

% ===================================================================
\section{Analytical benchmark: step-edge response}
% ===================================================================

As an independent check of the numerical implementation, the frequency
shift for a surface step edge ($\Delta z = 200$\,nm at $x = 0$) is
computed in closed form from the Poisson integral.  For
$z_s(x) = \Delta z\,\Theta(x)$ (Heaviside step), the second derivative
of the propagated potential yields:
\begin{equation}
\Delta\omega_x(x) = -\frac{e E_{\rm bg}\Delta z}{2m\omega_x}\,
    \frac{2}{\pi}\,\frac{x h}{(x^2 + h^2)^2},
\label{eq:step_solution}
\end{equation}
with extrema at $x = \pm h/\sqrt{3}$ and a full width at half maximum of
${\approx}0.34h$.

\textbf{Point spread function.}
The PSF is the derivative of the edge spread function,
$\mathrm{PSF}(x) = d\,\mathrm{ESF}/dx$, and follows in closed form:
\begin{equation}
\mathrm{PSF}(x) = -\frac{e E_{\rm bg}\Delta z}{2m\omega_x}\,
    \frac{2h}{\pi}\,\frac{h^2 - 3x^2}{(x^2+h^2)^3},
\label{eq:psf}
\end{equation}
with first zero at $x = h/\sqrt{3}$ (the ESF extremum), negative sidelobes
of relative amplitude ${\approx}0.09$, and asymptotic decay
$\mathrm{PSF}\sim 3x^{-4}$ --- the direct electrostatic analogue of an
optical PSF.  These quantities are distinct metrics, not competing definitions of a
single resolution: the ESF FWHM $0.34h$ is a characteristic width of the
edge response, the PSF first zero $x_0 = h/\sqrt{3}$ sets the central-lobe
scale, the Rayleigh-style two-point criterion based on the first zero
gives $\Delta x_{\rm res}=2h/\sqrt3\approx1.15h$, and the $50\%$ MTF
cutoff is $\lambda_{\min}\approx1.51h$; each is reported for the role it
plays in the main text.  Eq.~(\ref{eq:psf}) was verified against numerical
differentiation of the ESF to $4\times10^{-5}$ relative accuracy.

\textbf{Gaussian-ridge benchmark (fully closed form).}
For a 1D Gaussian ridge $z_s(x)=A\exp(-x^2/2\sigma^2)$, the half-space
Poisson integral evaluates in closed form:
\begin{equation}
\Phi^{(1)}(x,h) = E_{\rm bg}A\,
    \operatorname{Re}\bigl[e^{u^2}\operatorname{erfc}(u)\bigr],
\qquad
u = \frac{h - ix}{\sqrt{2}\,\sigma},
\label{eq:gauss_ridge}
\end{equation}
(the product $e^{u^2}\operatorname{erfc}(u)$, the scaled complementary
error function, is computed stably as $\operatorname{erfcx}(u)$), and
$\Delta\omega_x = -(e/2m\omega_x)\,\partial_x^2\Phi^{(1)}$ follows by
differentiating the recurrence
$d[e^{u^2}\operatorname{erfc}(u)]/du = 2u\,e^{u^2}\operatorname{erfc}(u)
-2/\sqrt\pi$.  For $A=100$\,nm, $\sigma=5$\,\textmu m, $h=40$\,\textmu m
the peak shift is $|\Delta\omega_x|/2\pi = 3488.2$\,Hz;
Eq.~(\ref{eq:gauss_ridge}) agrees with direct quadrature of the Poisson
integral to $10^{-15}$ relative accuracy and with the FFT forward model
to $6\times10^{-6}$ (peak-to-peak), providing an independent,
implementation-free check of the 2D code in the ridge limit.

Figure~\ref{fig:step} compares the analytical result
to the FFT-based numerical computation; the two agree to relative RMS
error $2.3\times10^{-3}$ (maximum $3.9\times10^{-3}$) over the edge
region $h/\sqrt3<|x|<3h$, the band $|x|<h/\sqrt3$ being excluded because
the analytic signal crosses zero there and the relative metric diverges;
the residual is the discretisation of the discontinuity and the periodic-wrap
baseline of the FFT domain.  Equation~(\ref{eq:step_solution}) provides the
\textbf{edge spread function} of the ion probe, confirming that the
lateral resolution is fundamentally limited to ${\sim}h$.

\begin{figure}[htb]
\centering
\includegraphics[width=\textwidth]{figV8_step_edge.pdf}
\caption{Analytical step-edge benchmark.
(a) Full step-edge response: $\Delta\omega_x(x)$ for a 200\,nm step.
The extrema occur at $x = \pm h/\sqrt{3}$, with FWHM ${\approx}0.34h$.
(b) Zoom at the edge, with key features annotated.
(c) Relative error on the masked domain $h/\sqrt3<|x|<3h$; the
zero-crossing band (shaded) is excluded because the relative metric
diverges there.}
\label{fig:step}
\end{figure}

\textbf{3D analytical benchmark: Gaussian bump.}
As a fully three-dimensional test of the ITF, consider a Gaussian surface
bump $z_s(x,y) = A\exp[-(x^2+y^2)/2\sigma^2]$.  The frequency shift can be
evaluated semi-analytically.  In Fourier space,
$\tilde{z}_s(k) = 2\pi A\sigma^2 e^{-\sigma^2 k^2/2}$.  The propagated
signal is $\Delta\tilde{\omega}_x = H_x \tilde{z}_s$, and the real-space
$\Delta\omega_x(x,y)$ is obtained by inverse FFT.  For $A = 100\,\textup{nm}$,
$\sigma = 5\,\mu\textup{m}$, $h = 40\,\mu\textup{m}$, the peak frequency shift
is $|\Delta\omega_x|/2\pi \approx 12\,\textup{Hz}$ --- detectable above the
$5\,\textup{Hz}$ noise floor with $N_{\rm rep} \sim 10$ repetitions.
This benchmark confirms that the ITF correctly handles fully 2D surface
features with azimuthal curvature, complementing the 1D step-edge test
and the exact Rayleigh cross-validation.  The FFT-based forward model reproduces
the semi-analytical result to within $10^{-12}$ relative error.

\textbf{Grid-convergence study.}
The numerical forward model uses a $128 \times 128$ grid with
$\Delta x = 0.5\,\mu\textup{m}$ (field of view $L = 64\,\mu\textup{m}$).
The smallest resolved wavelength at $h = 40\,\mu\textup{m}$ is
$\lambda_{\min} \approx 60\,\mu\textup{m}$, corresponding to
$\sim 120$ grid points per wavelength --- well within the Nyquist
criterion and far above the minimum ${\sim}10$ points per wavelength
required for $<1\%$ FFT error.  Varying $N$ from $64$ to $512$ changes
the reconstructed Pearson correlation by $<0.1\%$ for the synthetic test
case, confirming convergence.  The dominant numerical error source is not
the grid resolution but the periodic boundary condition implicit in the
FFT, which introduces $\mathcal{O}(e^{-L/h})$ edge effects from the
wrap-around of the $e^{-kh}$ filter.  For $L = 64\,\mu\textup{m}$ and
$h = 40\,\mu\textup{m}$, this error is ${\sim}2 \times 10^{-3}$ of the
peak signal, negligible for the present analysis.  A third independent
numerical validation is provided by a boundary-element (single-layer
collocation) solver with the periodic Green's function, which solves the
corrugated-boundary problem without any use of the ITF and reproduces the
exact Rayleigh result to ${\leq}0.6\%$ across the amplitude range
(Fig.~\ref{fig:bem_cross}).

% ===================================================================
\section{Calibration and traceability considerations}
% ===================================================================

A five-step calibration procedure establishes SI traceability:
(1)~trap field characterisation by measuring secular frequencies versus
applied dc voltages to extract $E_{\rm bg}$;
(2)~height calibration via electrical capacitance measurement and stage
interferometry;
(3)~ITF calibration using a PTB/NIST-certified sinusoidal grating with known
$A_{\rm ref}$, $\lambda_{\rm ref}$, fitting the measured $|\Delta\omega_x|$
to the ITF of the main text;
(4)~charge background characterisation by a reference scan at
$h_{\rm ref} \gg h_{\rm meas}$ where the topographic signal is negligible;
and (5)~verification against a certified step-height standard.
Traceability is to the SI metre (calibrated grating), SI second
(GPS-disciplined clock), and SI volt (Josephson voltage standard).

The traceability chain is:
\[
\textup{SI metre} \to \textup{laser interferometer} \to \textup{certified grating}
\to \textup{ITF calibration} \to z_s(x,y),
\]
\[
\textup{SI second} \to \textup{GPS-disciplined Rb clock} \to
\textup{rf source} \to \omega_x \to \Delta\omega_x,
\]
\[
\textup{SI volt} \to \textup{Josephson array} \to \textup{dc voltage source}
\to E_{\rm bg} \to \textup{surface potential scale}.
\]
The calibration uncertainty budget at each step is:
step~(1), $u(E_{\rm bg})/E_{\rm bg} \sim 10^{-3}$ (limited by dc voltage
source stability); step~(2), $u(h) \sim 100\,\textup{nm}$ (laser interferometry
+ capacitance); step~(3), $u(H_x)/H_x \sim 10^{-2}$ (grating certification
+ frequency measurement); step~(4), $u(\phi_{\rm charge}) \sim 1\,\textup{mV}$
(residual after UV cleaning); step~(5), $u(z_{\rm ref}) \sim 1\,\textup{nm}$
(certified step height).
The combined calibration uncertainty propagates into the measurement
uncertainty budget of the main text through the sensitivity
coefficients $c_i = |\partial z/\partial x_i|$.

\begin{table}[htb]
\centering
\caption{Component-wise degrees of freedom for the Welch--Satterthwaite
calculation of Sec.~3.14 of the main text.}
\label{tab:nu}
\small
\begin{tabular}{@{}lrr@{}r@{}}
\hline
Source & $u_i(z)$ [nm] & $\nu_i$ \\
\hline
Surface charge $\sigma$ & 319 & $\infty$ (rectangular) \\
ITF calibration (grating certificate) & 1.0 & 10 \\
Ion height $h$ & 2.0 & 30 (combined measurements) \\
Background field $E_{\rm bg}$ & 0.10 & 50 (calibration series) \\
Type A (frequency measurement) & 0.07 & $2(N_{\rm rep}-1)\approx1.1\times10^4$ \\
Secular frequency $\omega_x$ & $<10^{-3}$ & 20 \\
Ion mass $m$ & $<10^{-6}$ & $\infty$ (physical constant) \\
\hline
\end{tabular}
\end{table}

With $u_c = 319$\,nm, the Welch--Satterthwaite formula gives
$\nu_{\rm eff} = u_c^4\,/\sum_i u_i^4/\nu_i \approx 1.6\times10^{10}$,
dominated by the finite-$\nu$ entries of the table --- the Gaussian
coverage factor $k=2$ of the main text is fully justified.  The component
values $\nu_i$ are judgement-based, but the conclusion
$\nu_{\rm eff}\gg100$ is insensitive to them: it holds for any
$\nu_i\ge5$.

% ===================================================================
\section{Symbols and constants}
% ===================================================================

\begin{center}
\small
\begin{tabular}{@{}lll@{}}
\hline
Symbol & Meaning & Value (nominal) \\
\hline
$m$ & ion mass (${}^{40}$Ca$^+$) & $40$\,amu \\
$e$ & elementary charge & $1.602\times10^{-19}$\,C \\
$\omega_x,\omega_y$ & secular frequencies & $2\pi\times1$\,MHz \\
$E_{\rm bg}$ & background surface field & $10^4$\,V/m \\
$h$ & ion--surface standoff & $5$--$200\,\mu\textup{m}$ \\
$z_s(x,y)$ & surface height profile & \\
$\Phi^{(1)}$ & first-order boundary perturbation of the potential & \\
$\sigma(x,y)$ & surface charge density & $\lesssim10^{-3}$\,C/m$^2$ \\
$C = eE_{\rm bg}/2m\omega_x$ & ITF prefactor & $1.92\times10^3$\,m/s \\
$H_x(k_x,k_y;h)$ & instrument transfer function & $-Ck_x^2e^{-kh}$ \\
$\eta(k,h)$ & self-image correction & $\le0.15\%$ at $h=40\,\mu\textup{m}$ \\
$F(k)$ & dielectric-film factor & $1+kd(\varepsilon_r-1)/\varepsilon_r$ \\
$k_{\rm opt},\lambda_{\rm opt}$ & ITF peak & $2/h,\;\pi h$ \\
$\lambda_{\min}^{\rm(MTF)}$ & 50\% MTF short-wavelength cutoff & $1.51h$ \\
$\Delta k_{\rm FWHM}$ & ITF bandwidth & $3.395/h$ \\
$Q$ & passband quality factor & $0.589$ \\
$\delta\omega/2\pi$ & frequency noise floor & $5$\,Hz (conservative) \\
$\delta\omega_{\rm QPN}/2\pi$ & quantum projection noise floor & $0.7$\,Hz \\
$T_2$ & motional coherence time & $10$\,ms \\
$x_{\rm zp}$ & zero-point motion & $13$\,nm \\
$N,L,\Delta x$ & grid size, FOV, pixel & $128$, $64\,\mu\textup{m}$, $0.5\,\mu\textup{m}$ \\
$\sigma_n$ & singular values & main text \\
$r_{\rm eff}$ & effective rank & main text \\
$u_c(z), U$ & combined and expanded uncertainty & $319$\,nm,\;$638$\,nm \\
\hline
\end{tabular}
\end{center}

% ===================================================================
\section{Experimental roadmap}
% ===================================================================

The proposed implementation can be constructed from established
trapped-ion, precision-positioning and surface-characterisation
technologies, although their integration at the required ion--surface
separation presents significant experimental challenges.  The following
phased roadmap assumes existing surface-electrode Paul trap
infrastructure (cf.\ S{\"a}gesser et al., 2026):

\begin{enumerate}
\item \textbf{Phase~1 --- Static demonstration (0--12 months):}
      Position a ground-state-cooled ${}^{40}$Ca$^+$ at fixed height
      $h \approx 40\,\mu\textup{m}$ above a certified sinusoidal grating
      ($A = 50$\,nm, $\lambda = 20\,\mu\textup{m}$).  Measure
      $\Delta\omega_x$ via Ramsey interferometry with $\tau = 1$\,ms
      probe time, $N_{\rm rep} = 10^3$ repetitions.  Confirm the
      $k^2 e^{-kh}$ scaling by varying $h$ through the dc electrode
      voltages.  \textit{Expected signal:} the ITF predicts
      $|\Delta\omega_x|/2\pi = 5.3$\,Hz at $h=40$\,\textmu m --- at the
      conservative 5\,Hz floor, so the scaling test should start at
      $h\le30$\,\textmu m ($122$\,Hz at $h=30$\,\textmu m;
      $2.8$\,kHz at $h=20$\,\textmu m; $65$\,kHz at $h=10$\,\textmu m),
      where the single-pixel SNR exceeds $10$ for $N_{\rm rep}\sim10^2$.
      \textit{Key risk:} rf stability at the $10^{-6}$ level;
      \textit{mitigation:} stabilise rf source with a GPS-disciplined
      oscillator.

\item \textbf{Phase~2 --- 1D scanning (12--24 months):}
      Scan the ion along one axis ($x$) above a structured test surface
      with known topography.  Reconstruct the 1D profile via Tikhonov
      deconvolution and compare with AFM reference.  \textit{Key risk:}
      micromotion-induced position dependence of $\Delta\omega_x$;
      \textit{mitigation:} map micromotion via rf--photon correlation at
      each scan position and subtract the systematic.

\item \textbf{Phase~3 --- 2D scanning and full inversion (24--36 months):}
      Perform a full $128 \times 128$ pixel scan, apply 2D Tikhonov
      inversion, and compare the reconstructed $z_s(x,y)$ with AFM.
      Demonstrate the resolution--height trade-off (vary $h$ from
      $5$\,\textmu m to $100$\,\textmu m).  \textit{Key risk:}
      scan duration ($\sim 4.5$\,hours at 1\,s/pixel);
      \textit{mitigation:} implement real-time Bayesian adaptive scanning.

\item \textbf{Phase~4 --- Inter-laboratory comparison (36--60 months):}
      Two independent laboratories repeat the measurement on the same
      certified test surface.  Compare reconstructed $z_s(x,y)$ maps
      and uncertainty budgets.  Submit results to ISO/TC~213 as the
      basis for a new ISO~25178-6 subclass.  \textit{Key risk:}
      reproducibility of surface charge conditions; \textit{mitigation:}
      standardised UV-cleaning protocol before each measurement.
\end{enumerate}

\begin{table}[htb]
\centering
\caption{Phase-by-phase feasibility summary under the consistent noise model
of Sec.~3 of the main text ($\dot{\bar n}=10^2\,(h/40\,\mu\textup{m})^{-3.7}$
quanta/s): measurement duration, single-pixel signal-to-noise ratio
against the $5$\,Hz floor, and the dominant noise branch.  The heating
coherence times $T_2=1/\dot{\bar n}$ are $10$\,ms ($40$\,\textmu m),
$3.4$\,ms ($30$\,\textmu m), $0.77$\,ms ($20$\,\textmu m) and
$0.06$\,ms ($10$\,\textmu m).}
\label{tab:roadmap}
\small
\setlength{\tabcolsep}{4pt}
\begin{tabular}{@{}lp{3.6cm}cp{2.4cm}p{3.0cm}@{}}
\hline
Phase & Measurement & Duration & SNR (5\,Hz floor) & Dominant noise \\
\hline
1 & 1 point, $N_{\rm rep}=10^3$, $\tau=1$\,ms & 1\,s & 24 (122\,Hz at $h=30$\,\textmu m) & technical \\
2 & 1D scan, ${\sim}10^2$ points & ${\sim}10^2$\,s & $5.6\times10^{2}$ at $h=20$\,\textmu m & technical \\
3 & $128\times128$ at 1\,s/pixel & 4.55\,h & --- & technical ($h\ge40$\,\textmu m); heating ($h\lesssim5$\,\textmu m) \\
\hline
\end{tabular}
\end{table}

The total cost of the required apparatus (existing ion trap, laser
systems, and vacuum chamber) is dominated by the UHV system
(${\sim}\$200$\,k) and the laser setup (${\sim}\$150$\,k); the trap
chip and electronics add ${\sim}\$30$\,k.  This is comparable to a
research-grade AFM and substantially less than a cryogenic STM.

% ===================================================================
\section{Forward-model derivations}
% ===================================================================

\subsection{Second-order boundary perturbation}

$\Phi^{(2)}$ solves $\nabla^2\Phi^{(2)}=0$ in $z>0$ with Dirichlet datum
$\Phi^{(2)}(x,y,0)=-z_s\,\partial_z\Phi^{(1)}|_{z=0}$ (the quadratic
Taylor term of $\Phi^{(0)}$ vanishes because the background field is
uniform, $\partial_{zz}^2\Phi^{(0)}=0$).  From
$\widetilde{\partial_z\Phi^{(1)}}(\mathbf{k},0)
=-|\mathbf{k}|\,E_{\rm bg}\,\tilde z_s(\mathbf{k})$ and the convolution
theorem, the datum transform is
\begin{equation*}
\widetilde{\Phi^{(2)}}(\mathbf{k},0)=E_{\rm bg}\int\tilde
z_s(\mathbf{k}')|\mathbf{k}-\mathbf{k}'|\tilde
z_s(\mathbf{k}-\mathbf{k}')d^2\mathbf{k}'/(2\pi)^2,
\end{equation*}
and Poisson propagation multiplies by $e^{-kz}$.  The result was verified
against direct quadrature of the product $-z_s\,\partial_z\Phi^{(1)}|_{z=0}$
to machine precision, and its second harmonic at the ion reproduces the
$c_2$ coefficient of the exact Rayleigh expansion.

\subsection{Geometric majorant for the perturbation series}

Each normal derivative $\partial_z^{n+1-m}\Phi^{(m)}|_{z=0}$ acts as
multiplication by at most $k_{\max}^{n+1-m}$ in Fourier space, so the
recursion of the main text yields the majorant sequence
\begin{equation*}
a_{n+1}=\sum_{m=1}^{n}(z_{\max}k_{\max})^{n+1-m}\frac{a_m}{(n+1-m)!}
\end{equation*}
with $a_1=\|\Phi^{(1)}(\cdot,0)\|_{L^\infty}$.  Its generating function
$A(x)=\sum_{n\ge1}a_nx^n$ satisfies
\begin{equation*}
A(x)=a_1x+A(x)(e^{z_{\max}k_{\max}x}-1),
\end{equation*}
hence $A(x)=a_1x/(2-e^{z_{\max}k_{\max}x})$, whose radius of convergence
is the geometric bound $(\ln 2)/(z_{\max}k_{\max})$.

\subsection{Self-image (ion-induced) correction}

Applying the boundary perturbation of the main text to the charge--image
pair potential $\Phi_{\rm pair}=(e/4\pi\varepsilon_0)(r_1^{-1}-r_2^{-1})$,
whose normal derivative at the mean plane is
$\partial_z\Phi_{\rm pair}|_{z=0}=(e h/2\pi\varepsilon_0)(\rho^2+h^2)^{-3/2}$
with $\rho=|\mathbf{r}-\mathbf{r}_0|$, gives the induced boundary datum
$\psi=-z_s\,\partial_z\Phi_{\rm pair}|_{z=0}
= -z_s\,(e h/2\pi\varepsilon_0)(\rho^2+h^2)^{-3/2}$.  Its Fourier transform
follows from the Hankel transform
$\mathcal{F}[(\rho^2+h^2)^{-3/2}](q)=2\pi e^{-qh}/h$ as the convolution
\begin{equation*}
\tilde{\psi}(k)=-(e/4\pi^2\varepsilon_0)\int\tilde{z}_s(k')e^{-|k-k'|h}\,d^2k',
\end{equation*}
and Poisson propagation to the ion multiplies by $e^{-kh}$.  Evaluating
this convolution at the ion for a Fourier-mode surface and summing the
Hankel series yields the closed local form of the self-image lemma in the
main text.

% ===================================================================
\section{Quantum resource allocation and wave-packet smearing}
% ===================================================================

\subsection{Water-filling across modes}
Because the per-mode bound
$\delta z_k = 1/(|H_x(k)|\sqrt{\nu\mathcal{F}_{\rm ion}})$ varies by orders
of magnitude across the passband, a fixed quantum resource budget should
be allocated where it buys the most information.  With a total squeezing
budget $\sum_k r_k \le R_{\rm tot}$ and per-mode variance
$v_k e^{-2r_k}$ ($v_k \propto 1/|H_x(k)|^2$), the constrained minimisation
$\min_{r_k}\sum_k w_k v_k e^{-2r_k}$ is solved by the water-filling
prescription $r_k = \tfrac14\bigl[\ln v_k - \ln\mu\bigr]_+$ with the
Lagrange multiplier $\mu$ fixed by $\sum_k r_k = R_{\rm tot}$ --- the
quantum analogue of the classical channel water-filling of the information
capacity in the main text.  On the $h=40$\,\textmu m passband
(eight modes above the discrepancy-principle $\lambda$), a budget
$R_{\rm tot}=1.5$ ($13$\,dB) spent uniformly gives an average mode
variance $0.25$ (normalised units); the water-filling allocation
concentrates it on the two highest-$|H_x|$ modes and reduces the average
variance to $0.056$ --- a $\times 4.5$ information gain at equal
resources, verified numerically.

\subsection{Wave-packet smearing of the ITF}
The finite extent of the ion's wave packet convolves the measurement:
a thermal state with mean occupation $\langle n\rangle$ samples
$\Delta\omega_x(x)$ with Gaussian weight of variance
$\sigma_x^2 = x_{\rm zp}^2(2\langle n\rangle+1)$, so the effective
transfer function becomes
$H_x^{\rm(eff)}(k) = H_x(k)\,e^{-k^2\sigma_x^2/2}$ --- quantum smearing
adds its own Gaussian low-pass factor.  For ground-state cooling
($\sigma_x=x_{\rm zp}=13$\,nm), $k^2\sigma_x^2/2\le10^{-6}$ over the
accessible band and the correction is negligible; for a warm ion
($\langle n\rangle=10^3$, $\sigma_x\approx580$\,nm), the attenuation
reaches $0.7\%$ at $k=2\times10^5$\,rad/m and $4\%$ at the band edge ---
a systematic to be included in the error budget at elevated temperatures.
Because the smearing factor enters multiplicatively in the transfer
function, the quantum Fisher information acquires its square:
$\mathcal F_A^{\rm(eff)}(k)=\mathcal F_A(k)\,e^{-k^2\sigma_x^2}$, i.e.~the
quoted $0.7\%$ and $4\%$ attenuations reduce the QFI by $1.4\%$ and
$8\%$, respectively --- a temperature-dependent correction that enters the
quantum-enhancement claims verbatim.

% ===================================================================
\section{Stability and operating regime}
% ===================================================================

As the surface approaches the ion or the roughness amplitude increases, the
surface-induced force gradient eventually overcomes the trapping potential.
The stability criterion is $|\Delta\omega_x| < \omega_x/2$: beyond this
threshold, the secular frequency shift exceeds half the unperturbed value;
the secular motion then becomes unstable and the ion is ejected from
the trap.  Figure~\ref{fig:S_stability} maps the
stable and unstable regions in $(h, A)$ space for a sinusoidal surface with
$\lambda = 40$\,\textmu m.  The critical amplitude scales as
$A_{\rm crit}(h) \propto e^{+2\pi h/\lambda}$ for a surface with fixed
spatial period $\lambda$, growing rapidly with ion height.
(For the most-dangerous wavenumber $k = k_{\rm opt} = 2/h$,
$A_{\rm crit} \propto h^2$ with no exponential dependence.)

Three operating regimes follow: perturbation breakdown
($h \lesssim 5$\,\textmu m, $A \gtrsim 150$\,nm, where the linear ITF
fails), nominal sensing ($5$--$200$\,\textmu m, signal
$10^2$--$10^5$\,Hz), and signal-limited ($h \gtrsim 200$\,\textmu m,
below the $5$\,Hz floor).  At the nominal point ($h = 40$\,\textmu m,
$A = 100$\,nm) the stability margin exceeds two orders of magnitude.

\begin{figure}[htb]
\centering
\includegraphics[width=\textwidth]{figR1_stability.pdf}
\caption{Mathieu stability and QESPM operating regime.
(a) Standard $(a,q)$ stability diagram with three representative trap
operating points.  (b) Stability boundary in ion height--surface amplitude
$(h,A)$ space for a sinusoidal surface with $\lambda = 40$\,\textmu m.
The red curve marks $|\Delta\omega_x| = \omega_x/2$; the green and red
shaded regions indicate stable and unstable operation, respectively.}
\label{fig:S_stability}
\end{figure}

% ===================================================================
\section{Response to broadband surface roughness}
% ===================================================================

Real surfaces possess a continuous roughness spectrum characterised by
ISO~25178 parameters such as $R_a$.  Random profiles with controlled
$R_a$ (power-law spectrum $S(k) \propto [1+(k\ell_c)^2]^{-\beta}$,
$\ell_c = 5$\,\textmu m, $\beta = 1.5$) are propagated through the ITF
and the RMS frequency shift $\sigma_{\Delta\omega}$ is computed
(Fig.~\ref{fig:S_roughness}).

The response is \textbf{linear in $R_a$} for $R_a \ll h$, with a
roughness sensitivity of $68$\,Hz/nm at $h = 40$\,\textmu m rising to
$4.3\times10^3$\,Hz/nm at $h = 5$\,\textmu m: QESPM acts as a
\textbf{calibratable non-contact roughness gauge} linking the measured
shift to standardised texture parameters.  Linearity breaks down as
$R_a$ approaches $h/10$, where the boundary perturbation expansion fails.

\textbf{Exact linearity for stationary Gaussian surfaces.}
Let $z_s$ be a zero-mean stationary Gaussian random profile (1D scan) with
power spectral density $S(k)$, so that
$R_a^2 = (2/\pi)\langle z_s^2\rangle$.  Under the linear ITF,
Parseval's theorem gives the exact variance
\begin{equation}
\sigma_{\Delta\omega}^2 = \frac{1}{2\pi}\int_{-\infty}^{\infty}
    |H_x(k;h)|^2\,S(k)\,dk,
\label{eq:roughness_variance}
\end{equation}
where $H_x=-Ck^2 e^{-kh}$; because $S(k)\propto R_a^2$ for a fixed
spectral shape, $\sigma_{\Delta\omega}\propto R_a$ \emph{exactly} --- the
linear scaling is an identity of the linear model, verified to machine
precision against the FFT implementation and the Monte Carlo result
($68.0\pm2.5$\,Hz/nm at $h=40$\,\textmu m).  The onset of nonlinearity at
$R_a\sim h/10$ therefore signals \emph{model} breakdown (Sec.~S4), not
statistics.

\begin{figure}[htb]
\centering
\includegraphics[width=\textwidth]{figR2_roughness_response.pdf}
\caption{QESPM response to broadband surface roughness.
(a) RMS secular frequency shift $\sigma_{\Delta\omega}/2\pi$ as a function
of $R_a$ for ion heights $h = 5$--$80$\,\textmu m.  The linear scaling
$\sigma_{\Delta\omega} \propto R_a$ holds for $R_a \ll h$.  (b) Normalised
sensitivity $\sigma_{\Delta\omega}/R_a$, showing constant response in the
linear regime.}
\label{fig:S_roughness}
\end{figure}

% ===================================================================
\section{Sequential and adaptive measurement strategies}
% ===================================================================

For scanning applications, adaptive protocols can improve throughput.
A two-pass strategy is natural: (i)~a coarse scan at large $h$
($h \sim 80$\,\textmu m, $\Delta x \sim 10$\,\textmu m) identifies
regions of high topographic variance; (ii)~a fine scan at small $h$
($h \sim 5$--$20$\,\textmu m) resolves those regions at high resolution.
The Bayesian analogue updates the posterior over $z_s(x,y)$ sequentially
as each pixel is acquired, enabling real-time decisions about where to
scan next.  The greedy D-optimal selection is protected by a performance
guarantee: for Gaussian-process models the expected information gain is a
monotone submodular set function, and the greedy maximiser attains
$(1-e^{-1})\approx0.63$ of the optimal information gain
(Nemhauser--Wolsey-type bound; see Rasmussen and Williams, \emph{Gaussian
Processes for Machine Learning}).  For a fixed total measurement time
$T_{\rm tot}$, the optimal allocation between coarse and fine scans
follows from maximising the expected information gain, which depends on
the surface power spectrum.  A 1D demonstration with 17 partially
resolved Gaussian sources (spacing $5$\,\textmu m, below the PSF scale
$h=40$\,\textmu m) quantifies the gain (Fig.~\ref{fig:S_adaptive}): a
uniform 30-point acquisition achieves amplitude error $0.071$, while
greedy D-optimal point selection (maximising the expected information
gain $g^T P g$ at each step) with the same coarse budget plus 60 targeted
points achieves $0.028$ --- better than the exhaustive uniform 400-point
scan ($0.043$) at a quarter of the time --- and the posterior covariance
trace drops by a factor $4.6$.

\begin{figure}[htb]
\centering
\includegraphics[width=\textwidth]{figV12_adaptive.pdf}
\caption{Bayesian adaptive scanning (1D demonstration).
(a) True signal, noisy data, and the uniform-30 and adaptive (30+60 greedy
D-optimal) reconstructions.
(b) Posterior predictive standard deviation and the adaptively selected
refinement points.  Averaged over 30 noise realisations the amplitude
error drops from $0.071$ (uniform 30) to $0.028$ (adaptive), below the
exhaustive 400-point scan ($0.043$).}
\label{fig:S_adaptive}
\end{figure}

% ===================================================================
\section{Johnson--Nyquist noise: derivation details}
% ===================================================================

In the quasistatic limit ($\omega\ll c/d$, skin depth $\delta_s\gg d$), the
fluctuation--dissipation theorem applied to the reflected Green's function
of a dipole at height $d$ above a half-space of resistivity $\rho$, as
formulated for interfaces by Wylie and Sipe and applied to trapped ions by
Henkel et al., gives the $zz$-channel Green's function of the reflected
potential $G_\Phi(k;\omega)=(1/2\varepsilon_0 k)r_p(k,\omega)e^{-2kd}$ with
$\operatorname{Im} r_p = 2\varepsilon_0\omega/\sigma$ ($\sigma=1/\rho$,
Drude metal in the $\sigma\gg\varepsilon_0\omega$ limit), from which the
potential spectral density follows as
$S_\Phi(k;\omega;d)=(2k_BT/\omega)\operatorname{Im}G_\Phi
=2k_BT\rho\,k^{-1}e^{-2kd}$, and the field spectral density as
$S_E(k;\omega;d)=k^2S_\Phi=2k_BT\rho\,k\,e^{-2kd}$.  Integrating over $k$,
$S_E(\omega;d)=\int(k\,dk/2\pi)S_E(k)=k_BT\rho/(4\pi d^3)$, in agreement
with the FDT result $S_F=4k_BT\gamma$ for the Ohmic image drag
$\gamma=q^2\rho/(16\pi d^3)$ on a charge oscillating normally to the
surface.  The gradient observable is the $k_x^2$-weighted integral,
$S_{\partial E_x/\partial x}(\omega;d)=\int(k\,dk/2\pi)(k^2/2)S_E(k)
=3k_BT\rho/(8\pi d^5)$, giving the frequency-shift noise
$\delta\omega_J/2\pi=(e/4\pi m\omega_x)\sqrt{3k_BT\rho/(8\pi d^5\tau)}$
for a measurement of bandwidth $1/\tau$.

% ===================================================================
\section{Regularisation-theory proofs}
% ===================================================================

\textbf{Proof of the conditional (logarithmic) stability theorem.}
Split the frequency axis at the passband $B(\theta)=\{k:|H(k)|\ge\theta\}$.
On the passband, Parseval gives
$\|\widehat{z_1-z_2}\|_{L^2(B)}^2\le\delta^2/\theta^2$; on the complement,
the low-$k$ part is bounded through the embedding
$\|z\|_{L^\infty}\le C_s\|z\|_{H^s}$ and the high-$k$ part directly from
the source condition.  Combining the three pieces and balancing at
$\theta=\delta^{2/3}$ gives the announced bounds.

\textbf{Noise--bias decomposition of the Tikhonov estimate.}
The estimate obeys
$\|z_\lambda^\delta-z_s^\dagger\|_{L^2}\le\delta/(2\lambda)+
C_3E(\lambda/C)^{1/4}+Eh^s(\ln(Ch^2/\lambda))^{-s}(1+o(1))$: the first
term is the propagated noise (the Tikhonov filter satisfies
$\sup_k|H|/(|H|^2+\lambda^2)=1/(2\lambda)$) and the last two are the
regularisation bias from the two sides of the source set.

\textbf{Proof of the minimax lower bound.}
Two-hypothesis argument: the data of $z_\pm=\pm\varepsilon\cos(kx)$ differ
by $2\varepsilon|H(k)|$ and are indistinguishable from noise when
$\varepsilon\le\delta/|H(k)|$, while the source constraint requires
$\varepsilon(1+k^2)^{s/2}\le E$; balancing at
$k\approx h^{-1}\ln(1/\delta)$ gives the lower bound.

\end{document}
